\documentclass[11pt]{article}
\usepackage[margin=0.85in]{geometry}
\usepackage{amsmath,amssymb,amsthm,hyperref,array}
\newtheorem{theorem}{Theorem}
\setlength{\emergencystretch}{3em}
\newcommand{\mex}{\operatorname{mex}}
\newcommand{\dist}{\operatorname{dist}}
\title{A nonvacuous counterexample to the reverse-mex characterization}
\author{Myeongjin Kang}
\date{October 26, 2026}
\begin{document}\maketitle
\begin{abstract}
For subtraction Nim with moves $\{3,6,10\}$, we certify the exact Grundy sequence at every heap, the eventual reduction and positive-distance definitions from the source, and the complete universal condition in its proposed reverse-mex characterization. That condition holds nonvacuously, while reverse mex fails at heap five. Lean verifies the recurrence, uniqueness, every distance minimum, and the finite obstruction. The counterexample was already posted by Shivam Patel; no discovery or first-formalization claim is made.
\end{abstract}
\noindent\textit{Credit and license.} This AI-assisted note verifies Patel's existing counterexample. Original prose is CC BY 4.0 and original code is MIT.

For a finite subtraction set $S$, the source Grundy sequence is characterized by
\[
 g(n)=\mex\{g(n-s):s\in S,\ s\le n\},\qquad n\ge0,
\]
where mex is the least absent nonnegative integer. A heap is a $\mathcal P$-position when its Grundy value is zero. The reverse-mex property asks for
\[
 g(x)=\mex\{g(x+3),g(x+6),g(x+10)\}\quad\text{for every }x\ge0.
\]
\begin{theorem}
The game $S=\{3,6,10\}$ has eventual parameters $q=11,p=13$. It satisfies condition (a) of source Conjecture 1, with positive-distance minima, nonvacuously, but does not satisfy reverse mex.
\end{theorem}
\begin{proof}
The exact prefix is
\[
\begin{array}{c|rrrrrrrrrrrr}
n&0&1&2&3&4&5&6&7&8&9&10&11\\\hline
g(n)&0&0&0&1&1&1&2&2&2&0&3&3.
\end{array}
\]
Starting at $n=12$, repeat the block
\[
 (1,0,0,2,1,1,0,2,2,1,0,3,2).
\]
Here is an all-heap verification, rather than extrapolation of this displayed data. For $n<22$, direct substitution in the guarded mex recurrence verifies the displayed values. For $n\ge22$, all three predecessors are at least 12. Write $T(r)$ for the tail value at heap residue $r$ modulo 13. The complete tail check is
\[
\begin{array}{c|rrrrrrrrrrrrr}
r&0&1&2&3&4&5&6&7&8&9&10&11&12\\\hline
T(r)&0&0&2&1&1&0&2&2&1&0&3&2&1\\
T(r-3)&3&2&1&0&0&2&1&1&0&2&2&1&0\\
T(r-6)&2&1&0&3&2&1&0&0&2&1&1&0&2\\
T(r-10)&1&1&0&2&2&1&0&3&2&1&0&0&2
\end{array}
\]
with all arguments read modulo 13. In every column, the first row is the mex of the last three. This proves the defining recurrence at every heap, using its memory 10. Strong induction shows uniqueness: if two sequences satisfy the recurrence, then at heap $n$ all their legal predecessor values already agree. Consequently the displayed sequence is the actual source Grundy sequence, and $g(n+13)=g(n)$ for every $n\ge12$.

We now preserve the source's eventual reduction exactly. For an integer target $z$, define
\[
 R(z)=\begin{cases}z,&z\ge12,\\12+((z-12)\bmod13),&z<12,\end{cases}
\]
where the second branch uses the remainder in $\{0,\ldots,12\}$. If $z<12$, this is $z+13h$ for a positive integer $h$ and lies in $[12,24]$; if $z\ge12$, the target is left unchanged. This is Definition 12 with $q=11,p=13$. In both cases the reduced heap is in the tail and has residue $z$ modulo 13.

Distance is the least \emph{positive} integer $k$ with $g(R(w-3k))=0$, as in the source's positive-natural convention. A minimum exists for every $w$: for instance $k=9w+13$ gives $w-3k=-13(2w+3)$ and hence a zero-valued tail residue. Computing the first positive hit in each residue gives the complete table
\[
\begin{array}{c|rrrrrrrrrrrrr}
w\bmod13&0&1&2&3&4&5&6&7&8&9&10&11&12\\\hline
\dist(w)&4&3&2&1&1&3&2&2&1&3&3&2&1.
\end{array}
\]
Every table entry is a minimum: substitution checks the indicated hit and all smaller positive $k$. This table applies also to the prefix heaps because the reduction always enters the tail.

The prefix $\mathcal P$-positions are $0,1,2,9$. The tail $\mathcal P$ residues are $0,1,5,9$. Combining the full position classification with the distance table shows that the $\mathcal P$-positions with $\dist(w)=2m$, $m\ge2$, are precisely
\[
 w=0,13,26,39,\ldots,
\]
and their distance is four. At every such heap,
\[
 g(R(w+10-3))=T(7)=2\ne0,\qquad
 g(R(w+10-3\cdot3))=T(1)=0.
\]
Thus the universal source condition (a) holds. It is nonvacuous both at heap 0 and at the positive heap 13; no prefix position has been omitted.

Finally,
\[
 g(5)=1\ne0=\mex\{2,3,2\}
   =\mex\{g(8),g(11),g(15)\}.
\]
Reverse mex therefore fails at heap 5, although condition (a) holds nonvacuously. This contradicts the proposed equivalence.
\end{proof}
The module \texttt{ReverseMex.lean} defines mex by a genuine least-number construction and proves its membership/minimality characterization. The actual legal move set is the filtered set $S=\{3,6,10\}$ followed by the predecessor-value image. The module proves the bridge to the guarded recurrence, verifies the complete finite prefix and all 13 tail cases, and proves uniqueness for every source SG sequence. It also derives the normal-play losing-position recursion from the mex equation. The reduction specification proves both branches of Definition 12, including the positive integral lift. The distance theorem proves every hit and its positive minimality, followed by the full qualifying-position classification and universal condition. The final theorem is \texttt{ReverseMexCounterexample.source\_counterexample}.

The construction and mathematical counterexample were already posted by Shivam Patel. This certificate proves the single base game and does not claim the posted infinite-scaling strengthening. It does not refute the source's separate proved special family or its pass-game assertions. Source URLs, exact solution identifiers, hashes, and pinned Lean 4.34.1 reproduction instructions accompany the certificate.
\begin{thebibliography}{9}
\bibitem{source} U. Larsson, H. Manabe, and R. Miyadera, \emph{A Subtraction Nim with a Pass}, arXiv:2605.14321v1, Definitions 4, 8, 12--13 and Section 4.1, Conjecture 1. \url{https://arxiv.org/abs/2605.14321v1}.
\bibitem{post} S. Patel, posted counterexample to MathDB 374728. \url{https://mathdb.com/p/374728}.
\end{thebibliography}
\end{document}
